% Figure 1.1 : machines virtuelles et conteneurs, ce qui est dupliqué et ce qui est partagé.
\documentclass[tikz,border=4pt]{standalone}
\usepackage{quai}
\begin{document}
\begin{tikzpicture}[x=1cm,y=1cm,
    couche/.style={qbox, minimum height=0.78cm, font=\footnotesize, inner sep=3pt},
    app/.style={couche, draw=QVert, fill=QVertBg},
    lib/.style={couche, draw=QBleu, fill=QBleuBg, align=center, text width=2.9cm},
    noyau/.style={couche, draw=QOrange, fill=QOrangeBg},
    base/.style={couche, draw=QGris, fill=QGrisBg}]

  % ---------------- machines virtuelles ----------------
  \node[qtitre] at (0,7.6) {Machines virtuelles};
  \node[base, minimum width=7.2cm] at (3.6,0.4) {matériel (processeurs, mémoire, disque)};
  \node[base, minimum width=7.2cm] at (3.6,1.3) {système hôte et hyperviseur (KVM, VirtualBox...)};
  \foreach \x/\a/\l in {1.75/application A/Python 3.14, 5.45/application B/Python 3.9} {
    \draw[qcadre, dashed] (\x-1.7,1.85) rectangle (\x+1.7,6.85);
    \node[qlbl, anchor=north west] at (\x-1.6,6.8) {machine virtuelle};
    \node[noyau, minimum width=3.1cm] at (\x,2.4) {noyau invité};
    \node[lib, minimum width=3.1cm] at (\x,3.4) {système invité\\complet};
    \node[lib, minimum width=3.1cm] at (\x,4.55) {\l\\et ses bibliothèques};
    \node[app, minimum width=3.1cm] at (\x,5.6) {\a};
  }
  \node[qlblc=QOrange, align=center] at (3.6,-0.55) {chaque machine démarre son propre noyau\\et son propre système d'exploitation};

  % ---------------- conteneurs ----------------
  \begin{scope}[xshift=8.6cm]
    \node[qtitre] at (0,7.6) {Conteneurs};
    \node[base, minimum width=7.2cm] at (3.6,0.4) {matériel (processeurs, mémoire, disque)};
    \node[noyau, minimum width=7.2cm, line width=1.1pt] at (3.6,1.3) {un seul noyau Linux, partagé};
    \node[base, minimum width=7.2cm] at (3.6,2.2) {moteur de conteneurs (Docker, containerd, Podman)};
    \foreach \x/\a/\l in {1.75/application A/Python 3.14, 5.45/application B/Python 3.9} {
      \draw[qcadre, dashed] (\x-1.7,2.75) rectangle (\x+1.7,5.85);
      \node[qlbl, anchor=north west] at (\x-1.6,5.8) {conteneur};
      \node[lib, minimum width=3.1cm] at (\x,3.45) {\l\\et ses bibliothèques};
      \node[app, minimum width=3.1cm] at (\x,4.6) {\a};
    }
    \node[qlblc=QOrange, align=center] at (3.6,-0.55) {chaque conteneur n'est qu'un processus du noyau\\hôte, qui voit ses propres fichiers};
  \end{scope}

  % légende
  \begin{scope}[shift={(2.2,-1.75)}]
    \foreach \x/\c/\cbg/\t in {0/QVert/QVertBg/votre application,
        4.1/QBleu/QBleuBg/ce que l'application emporte avec elle,
        10.2/QOrange/QOrangeBg/noyau} {
      \draw[draw=\c, fill=\cbg, line width=0.8pt] (\x,0) rectangle ++(0.45,0.3);
      \node[qlbl, anchor=west] at (\x+0.55,0.15) {\t};
    }
  \end{scope}
\end{tikzpicture}
\end{document}
